\begin{textsample}{POS Dim 3 – global\_south\_2024\_01 – Score 7.00 – global\_south\_2024\_01/105031556.txt}  \label{ex:f3_pos_381}
Israel described the allegations levelled by South Africa as hypocritical and said that one of the biggest cases ever to come before an international court reflected a world turned upside down.
Legal adviser to Israel ' s Foreign Ministry Tal Becker looks to his right inside the International Court of Justice ( ICJ ) as judges hear a request for emergency measures by South Africa to order Israel to stop its military actions in Gaza, in The Hague, Netherlands January 12, 2024. ( Reuters Photo ) Accused of committing genocide against Palestinians, Israel insisted at the United Nations ' highest court Friday that its war in Gaza was a legitimate defense of its people and said instead that Hamas was guilty of genocide.
Israel described the allegations levelled by South Africa as hypocritical and said that one of the biggest cases ever to come before an international court reflected a world turned upside down.
Israeli leaders defend their air and ground offensive in Gaza as a legitimate response to Hamas ' October 7 attack, when militants stormed through Israeli communities, killed some 1,200 people and took @ @ @ @ @ @ @ @ @ @ told a packed auditorium at the ornate Palace of Peace in The Hague that the country is fighting a"war it did not start and did not want.""In these circumstances, there can hardly be a charge more false and more malevolent than the allegation against Israel of genocide,"he added, noting that the horrible suffering of civilians in war was not enough to level a charge of genocide.
Advertisement South African \textbf{lawyers} asked the court Thursday to order an immediate halt to Israeli military operations in the besieged coastal territory that is home to 2.3 million Palestinians.
A decision on that request will probably take weeks, and the full case is likely to last years -- though it ' s unclear if Israel would follow any court orders.
On Friday, Israel focused on the brutality of the October 7 attacks, presenting chilling video and audio to a hushed audience to highlight what happened that day."They tortured children in front of parents and parents in front of children, burned @ @ @ @ @ @ @ @ @ @ mutilated scores of women, men and children,"Becker said.
South Africa ' s request, he said, amounts to an attempt to prevent Israel from defending against that assault.
Even when acting in self-defense, countries are required by international law to follow the rules of war, and the court must decide if Israel has.
Advertisement Israel often boycotts international tribunals and UN investigations, saying they are unfair and biased.
But this time, Israeli leaders have taken the rare step of sending a high-level legal team -- a sign of how seriously they regard the case and likely their fear that any court order to halt operations would be a major blow to the country ' s international standing.
Still, Becker dismissed the accusations as crude and attention-seeking."We live at a time when words are cheap in an age of social media and identity politics.
The temptation to reach for the most outrageous term to vilify and demonize has become, for many, irresistible,"he said. @ @ @ @ @ @ @ @ @ @."If there have been acts that may be characterized as genocidal, then they have been perpetrated against Israel,"Becker said.
Hamas has, he said, a"proudly declared agenda of annihilation, which is not a secret and is not in doubt."More than 23,000 people in Gaza have been killed during Israel ' s military campaign, according to the Health Ministry in the Hamas-run territory.
Nearly 85 per cent of Gaza ' s people have been driven out of their homes, a quarter of the enclave ' s residents face starvation, and much of northern Gaza has been reduced to rubble.
Advertisement South Africa says this amounts to genocide and is part of decades of Israeli oppression of Palestinians."The scale of destruction in Gaza, the targeting of family homes and civilians, the war being a war on children, all make clear that genocidal intent is both understood and has been put into practice.
The articulated intent is the destruction of Palestinian life,"said \textbf{lawyer} Tembeka Ngcukaitobi @ @ @ @ @ @ @ @ @ @ the fighting and Israel didn't comply, it could face UN sanctions, although those may be blocked by a veto from the United States, Israel ' s staunch ally.
The White House declined to comment on how it might respond if the court determines Israel committed genocide.
But National Security Council spokesperson John Kirby called the allegations"unfounded."The extraordinary case goes to the core of one of the world ' s most intractable conflicts -- and for the second day \textbf{protesters} \textbf{rallied} outside the court.
Pro-Israeli \textbf{demonstrators} set up a table near the court grounds for a Sabbath meal with empty seats commemorating the hostages still being held by Hamas.
Nearby, over 100 \textbf{pro-Palestinian} \textbf{protesters} waved flags and shouted \textbf{protests}.
The case also strikes at the heart of both Israel ' s and South Africa ' s national identities.
Israel was founded as a Jewish state in the wake of the Nazis ' slaughter of 6 million Jews during World War II.
South Africa ' s governing party, meanwhile, has long compared Israel ' s policies in @ @ @ @ @ @ @ @ @ @ the apartheid regime of white minority rule, which restricted most Black people to"\textbf{homelands}."The world court, which rules on disputes between nations, has never judged a country to be responsible for genocide.
The closest it came was in 2007, when it ruled that Serbia"violated the obligation to prevent genocide"in the July 1995 massacre by Bosnian Serb forces of more than 8,000 Muslim men and boys in the Bosnian enclave of Srebrenica.

% matched lemmas: demonstrator, homeland, lawyer, pro-palestinian, protest, protester, rally
\end{textsample}
